\documentclass{article}

\usepackage[english]{babel}
\usepackage[letterpaper,top=2cm,bottom=2cm,left=3cm,right=3cm]{geometry}
\usepackage{amsmath,amssymb}
\setlength{\parindent}{0pt}

\title{DSC 190/291: Learning Theory through Formal Proof and Proof Presentation\\Assignment 2}
\date{UCSD, Spring 2026}

\begin{document}
\maketitle

\section{Overview}

This assignment has two technical parts plus an AI usage report:
\begin{itemize}
\item Part A: unions of two intervals on the line.
\item Part B: quadratic threshold functions on the line.
\item Part C: AI usage report.
\end{itemize}

\section{Part A: Unions of Two Intervals on the Line (40 points)}

Let
$$
\mathcal{H}_2
=
\{x \mapsto 1[x \in I_1 \cup I_2] : I_1, I_2 \subseteq \mathbb{R} \text{ are intervals}\}.
$$
In other words, a hypothesis in $\mathcal{H}_2$ labels a point by $1$ iff it lies in a union of at most two intervals on the real line. Empty intervals are allowed, so this really means ``at most two.''

Throughout this part, fix ordered points
$$
x_1 < x_2 < \cdots < x_n.
$$

\begin{enumerate}
\item \textbf{Restriction patterns.}

Characterize exactly which binary labelings of $(x_1,\dots,x_n)$ can be realized by $\mathcal{H}_2$.
A good answer should not just list examples. It should identify the structural property that distinguishes realizable labelings from impossible ones.

\item \textbf{Exact growth function.}

Use your characterization to derive a closed-form formula for the growth function
$$
\Gamma_{\mathcal{H}_2}(n).
$$
You may first use small cases to guess the formula, but your final answer must be a proof.

\item \textbf{VC dimension and mistake bound.}

Determine the exact value of
$$
\operatorname{VCdim}(\mathcal{H}_2).
$$
Then use your exact growth formula to write down the Halving mistake bound on an $n$-point pool in the realizable online transductive setting.

Finally, compare your exact growth formula with the Sauer--Shelah bound obtained from your VC-dimension calculation. Are they equal in this case? What does this tell you about when a VC-dimension-based upper bound can be tight?

\item \textbf{Tightness of Sauer--Shelah.}

For arbitrary integers $n \ge d \ge 0$, give a concrete hypothesis class $\mathcal{H}$ on some domain $\mathcal{X}$ such that
$$
\operatorname{VCdim}(\mathcal{H}) = d
$$
and
$$
\Gamma_{\mathcal{H}}(n) = \sum_{k=0}^d \binom{n}{k}.
$$
What does your example show about the Sauer--Shelah bound?

\item \textbf{AI proof audit.}

An AI assistant claims:

\begin{quote}
On an ordered $n$-point sample, a labeling realized by $\mathcal{H}_2$ is determined by the places where the labels switch between $0$ and $1$. Since a union of at most two intervals can create at most four such switches, one just chooses up to four switch locations among the $n$ sample positions. Therefore
$$
\sum_{j=0}^4 \binom{n}{j}.
$$
Hence
$$
\Gamma_{\mathcal{H}_2}(n) = \sum_{j=0}^4 \binom{n}{j},
$$
and in particular
$$
\operatorname{VCdim}(\mathcal{H}_2) = 4.
$$
\end{quote}

A flawed argument may still arrive at a true conclusion; analyze the reasoning, not just the final claim.

Explain carefully what is incomplete or incorrect in this argument. Then replace it with a correct statement that is actually supported by your work above.
\end{enumerate}

\section{Part B: Quadratic Threshold Functions on the Line (45 points)}

Let
$$
\mathcal{H}_{\mathrm{quad}}
=
\{x \mapsto 1[ax^2 + bx + c \ge 0] : (a,b,c) \in \mathbb{R}^3,\ (a,b,c) \ne (0,0,0)\}.
$$
So a hypothesis in $\mathcal{H}_{\mathrm{quad}}$ is the indicator of the nonnegative region of a quadratic polynomial.

Throughout this part, again fix ordered points
$$
x_1 < x_2 < \cdots < x_n.
$$

\begin{enumerate}
\item \textbf{A general upper-bound trick.}

Prove the following statement.

If there exist an integer $D \ge 1$ and a transformation
$$
\phi : \mathcal{X} \to \mathbb{R}^D
$$
such that every hypothesis $h \in \mathcal{H}$ can be written in the form
$$
h(x) = 1[\langle w, \phi(x) \rangle \ge 0]
$$
for some vector $w \in \mathbb{R}^D$, then
$$
\operatorname{VCdim}(\mathcal{H}) \le D.
$$
You may use the Week 2 result that homogeneous halfspaces in $\mathbb{R}^D$ have VC dimension $D$.

\item \textbf{Apply the trick to quadratic thresholds.}

Find an explicit transformation
$$
\phi : \mathbb{R} \to \mathbb{R}^3
$$
such that every hypothesis in $\mathcal{H}_{\mathrm{quad}}$ can be written in the form
$$
h(x) = 1[\langle w, \phi(x) \rangle \ge 0]
$$
for some $w \in \mathbb{R}^3$.

Then use the previous item to prove an upper bound on
$$
\operatorname{VCdim}(\mathcal{H}_{\mathrm{quad}}).
$$

\item \textbf{Exact restriction patterns.}

Characterize exactly which binary labelings of $(x_1,\dots,x_n)$ are realizable by $\mathcal{H}_{\mathrm{quad}}$.
Your characterization should make clear how the algebraic structure of a quadratic polynomial controls the combinatorics of its restriction to an ordered finite set.

\item \textbf{Exact growth and exact VC dimension.}

Use your characterization to compute the exact growth function
$$
\Gamma_{\mathcal{H}_{\mathrm{quad}}}(n)
$$
and the exact VC dimension of $\mathcal{H}_{\mathrm{quad}}$.

Then compare your exact growth formula with the Sauer--Shelah bound obtained from your VC-dimension calculation. Is the VC-based upper bound tight here? Explain what your answer says about VC dimension as a summary of finite-pool richness.

\item \textbf{AI proof audit.}

An AI assistant claims:

\begin{quote}
A quadratic polynomial has at most two real roots, so on an ordered sample its labels can change from $0$ to $1$ or from $1$ to $0$ at most twice. Therefore one chooses up to two change-points among the $n-1$ gaps and gets
$$
\Gamma_{\mathcal{H}_{\mathrm{quad}}}(n)
=
\sum_{j=0}^2 \binom{n-1}{j}.
$$
In particular,
$$
\operatorname{VCdim}(\mathcal{H}_{\mathrm{quad}}) = 3.
$$
\end{quote}

A flawed argument may still arrive at a true conclusion; analyze the reasoning, not just the final claim.

Explain carefully what is incomplete or incorrect in this argument. Then replace it with a correct theorem that is genuinely justified by your work in this part.
\end{enumerate}

\section{Part C: AI Usage Report (15 points)}

Write a short report describing how you used AI in this assignment. Do not just list tools; explain what role AI played in your work and how you checked the result. Address:
\begin{enumerate}
\item Describe the parts of the assignment for which you used AI.
\item Describe concrete AI suggestions you accepted and explain why.
\item Describe concrete AI suggestions you rejected or substantially modified, and explain what was wrong, incomplete, or unhelpful about them.
\item Describe how you verified the correctness of what you submitted.
\item Describe concrete updates to your AI workflow that resulted from this assignment.
\end{enumerate}

If you did not use AI for some part of the assignment, say so explicitly.

\end{document}
